\section{The target game: Gomoku under Swap2}
\label{sec:gomoku}

\subsection{Rules}
\label{sec:rules}

The game is played on a $15\times15$ board. Two players, Black and White,
alternate placing one stone of their colour on an empty cell. The first player to
make five or more of their stones in an unbroken line horizontally, vertically,
or diagonally wins. A line of six or more -- an \emph{overline} -- counts as a
win; this is the \emph{freestyle} rule set, and it differs from the Renju rule
set, in which an overline is not a win for Black and several move types are
forbidden \citep{renju_rules}. The board has 225 cells, so a game that fills the
board without a line is a draw; with $\gamma = 1$ and outcomes in $\{-1, 0, +1\}$,
the setup is exactly the MDP specialization of Section~\ref{sec:mdp}.

Four decisions in the rules affect the implementation and the learning problem,
and each is worth stating while the rules are still in view.

Overlines counting as wins means the win test is ``five or more in a row'' rather
than ``exactly five'', which is a strictly simpler test: one that a bitmask
formulation expresses as a single comparison rather than as a comparison plus a
forbidden-neighbour check.

A draw is reachable, if rare, so the value head's target set is $\{-1,0,+1\}$ and
the last term of \eqref{eq:alphazero-loss} must handle $z = 0$ without special
treatment. Mean squared error does this naturally.

The rules are translation-invariant and invariant under the board's dihedral
group: the eight rotations and reflections of the square. This is what makes
symmetry augmentation legitimate and what makes the plane encoding and the policy
output transform-compatible.

Finally, the rules are \emph{forcing} at a structural level that Go is not. A
threat that would complete five in a row on the next move \emph{must} be answered,
so the defender's branching factor collapses to one at exactly the moments when
the game is decided. Section~\ref{sec:alphaepsilon} turns this structural fact
into a component.

\subsection{A solved game, and why we play it anyway}
\label{sec:solved}

In 1993 and 1994 Allis, van den Herik, and Huntjens proved that freestyle Gomoku
on $15\times15$ is a forced win for the first player
\citep{allis1996gomoku,allis1994thesis}. The method was threat-space search
combined with proof-number search: rather than search the full game tree, the
argument exploits the forcing structure described above, so that the attacker
considers only moves that create a threat and the defender only the forced
replies. This reduced a tree far beyond brute-force reach to something provable.

That the empty board is a first-player win does not trivialize the problem, for
three reasons.

First, competitive play does not start from the empty board. It uses an opening
protocol -- Swap, Swap2 -- or restriction rules, precisely because the bare game
is unbalanced. Under the Swap2 protocol the game has \emph{not} been solved, and
between the empty board and the solved opening lines the theoretical values of
most legal positions are unknown. A strong Swap2 agent operates on genuinely open
ground.

Second, a solver is not a learning agent. Threat-space search is a hand-crafted
exploit of one game's structure. AlphaZero is the opposite kind of object: a
general method that discovers structure. The scientific interest lies in the
generality -- the same code should play any two-player perfect-information board
game after the rules engine is exchanged.

Third, a solved game supplies free test data. Positions with known exact values
and forced-win sequences that can be verified mechanically are ideal test cases
for a search implementation, and they exist before any learning has taken place.
An agent that cannot win a won position is defective, and that can be established
with a test suite and no gradients.

\paragraph{The Swap2 protocol.} One player places three stones -- two of their
own colour and one of the opponent's -- in any order and any positions. The other
player then has three choices: take Black, take White, or place two more stones
(one of each colour) and pass the colour choice to the first player. The protocol
is a bidding mechanism: whoever believes the position favours them offers to take
the disadvantaged colour.

For an implementation, Swap2 has one important structural consequence: the first
few moves do not alternate between players. A board representation that stores
positions relative to the side to move -- the ``me/you'' view that
Section~\ref{sec:architecture} describes as an input encoding -- cannot express
the opening, because during the opening the same player places two stones in
succession. The rules engine must therefore store \emph{absolute} colours, and
the relative view is derived at encoding time. Separate the two concerns: the
engine stores colours, the encoder produces the player-relative planes.

\subsection{Why Gomoku suits the method}
\label{sec:gomoku-suits}

Three properties make Gomoku a favourable target, and they are the reason a
modest budget produces a non-embarrassing agent.

\emph{The horizon is short.} A game lasts roughly 30 to 70 moves. AlphaGo Zero
used $\tau = 1$ for the first 30 moves of games that lasted around 200 moves, so
that openings stayed diverse while the majority of the game was played greedily.
Scaled to a 60-move game the equivalent phase is 10 to 12 moves, and the same
argument applies with the same force. This is a scaling of a hyperparameter, not
a change of algorithm, and it is the least controversial of the transfers.

\emph{The tactics are shallow and recency-driven.} Influence in Go accumulates
over hundreds of moves; in Gomoku, one forcing sequence decides the game, and it
usually involves the last few moves. The practical consequence is that the
history planes of AlphaGo Zero's encoding are largely wasted here: eight time
steps of history cost channels and buy little, while the last two moves carry
most of the temporal information a player needs. Published AlphaZero-style
Gomoku agents use three planes; a small network of 32 filters and two residual
blocks reached human playing strength in two days on one 2018-era GPU
\citep{xie2018alphagomoku}. The architectural lesson is encouraging: two orders of
magnitude fewer parameters than AlphaGo Zero.

That paper also reports that a curriculum mattered -- imitation of synthetic
attack and defence moves first, then a rule-based mentor, then self-play -- which
is the second reason for the component in Section~\ref{sec:alphaepsilon}.

\emph{The state space is small enough to test exhaustively at the level that
matters.} Differential testing between a fast bitboard implementation and a naive
array implementation is cheap and decisive: ten thousand random games compared
ply by ply will find any disagreement between the two, and a disagreement is a
bug in one of them. This is the practical form of ``verify before you claim'' for
a rules engine, and it is available because Gomoku's rules are simple enough to
write twice in different styles.

\subsection{The alpha-epsilon position}
\label{sec:alphaepsilon}

A pure AlphaZero agent knows only the rules. The design described here keeps a
small amount of hand-written game knowledge, and the honest name for such a
system is not tabula rasa. What follows states precisely what the component may
and may not do, because the value of a small amount of hand-written knowledge
depends entirely on keeping it small and keeping it honest.

\paragraph{What it is.} Three predicates, each computed by bit operations over
the two players' stone sets:

\begin{itemize}
  \item \texttt{immediate\_wins(side)}: the empty cells where placing a stone
        completes five in a row for \code{side}.
  \item \texttt{forced\_blocks()}: the opponent's immediate wins, which are
        therefore the moves the side to move is compelled to play.
  \item \texttt{double\_threats(side)}: the cells after which \code{side} has at
        least two immediate wins, which is unanswerable -- an open four or a
        double four.
\end{itemize}

All three are instances of one primitive: place a stone hypothetically, and ask
whether a win exists. On a 225-cell board that is roughly 225 cheap tests per
call, which is inexpensive at leaf expansion and far too expensive inside a
selection loop.

\paragraph{What it is used for.} Four roles, in order of how much they matter.

As a \emph{deterministic stand-in evaluator}, it lets the whole search be built
and tested before a network exists. The mock evaluator assigns high priors and
near-certain values to immediate wins and forced blocks and mild shaping
elsewhere. The search can then be checked against tactical puzzles with known
solutions: a correct implementation given a sane prior must find a win in one at
50 simulations, must play the forced block, and must find short forced sequences.
This is the single most valuable use, and it is a testing use, not a playing use.

As a \emph{prior shaping term}, it gives a young network's near-uniform priors a
useful tilt early in training. The blend weight should decay over iterations:
hand-written knowledge is a scaffold, and a scaffold that never comes down limits
the structure that can be learned.

As a \emph{label generator}, it seeds a supervised phase with exact labels before
self-play data is abundant. Here the component must be tightened, and that is
what the next paragraph is about.

As a \emph{fast path} during self-play, an immediate win can be played without a
search, saving a full budget of simulations on a move whose value is already
known.

\paragraph{The prover, and the soundness rule.} The detection predicates say what
is true now. A bounded threat-space search says what is \emph{forced}: it
recursively searches the attacker's threat-creating moves and the defender's
forced replies, and returns a certified winning sequence when it finds one.
Because the defender's branching factor is approximately one, the search is
small, and a node-and-depth budget is enough to keep it at microsecond-to-
millisecond cost.

Two properties make such a prover safe to use as a label generator, and both must
be stated as policy rather than assumed.

\emph{Soundness over completeness.} Exhausting the budget must mean ``no proof
found'', never ``loss''. A missed win costs puzzle coverage; a wrong label
poisons the value head with false certainty, which is far more expensive, because
the network is the search's evaluator and a confidently wrong value propagates
into selection.

\emph{Every label is machine-verified.} A proof is emitted only after replaying
the line while enumerating all defender alternatives at each defender turn. The
replay is cheap precisely because the replies are forced. A proof object should
behave like a witness: it carries the line, and it is checked rather than
trusted.

A verified proof becomes training data with an exact value target $z = \pm1$,
which is not a bootstrap estimate. Such positions belong in an \emph{anchor set}:
a small fraction of every batch, permanently retained rather than subject to
replay-window eviction. Exact labels are the one kind of supervision whose value
does not decay as the network improves.

\paragraph{What it must not do.} The component must not encode strategy. It has
no notion of a good opening, no positional evaluation, and no search beyond the
bounded prover. Its purpose is to remove the first few thousand games of obvious
mistakes, and to make the search testable before training exists. Anything beyond
that is not epsilon; it is a solution to the game, and there is already one.

\subsection{Scaling the hyperparameters}
\label{sec:scaling}

Table~\ref{tab:scaling} transfers the published constants to Gomoku. Two entries
deserve comment.

The Dirichlet concentration follows AlphaZero's own rule, $\alpha \approx
10/(\text{typical legal moves})$, which yields $0.03$ at 361 Go intersections,
$0.15$ at roughly 80 shogi moves, and $0.3$ at roughly 35 chess moves
\citep{silver2018alphazero}. Gomoku's midgame offers roughly 100 to 150 legal
moves, so $10/125 \approx 0.08$; the rule gives a smaller $\alpha$ than Go's,
which is the correct direction, since a larger action set wants a spikier noise
vector to concentrate the exploration.

The simulation count is a pure cost--quality trade. AlphaZero used 800
simulations per move with a well-trained network. A weaker network benefits less
from deep search, because deep search over a poor evaluator mostly propagates the
evaluator's errors, while self-play throughput benefits directly from a smaller
budget: halving simulations halves the cost of every game. Four hundred
simulations per move is a reasonable starting point for a first working system,
to be raised once the network is strong enough to reward depth.

\begin{table}[t]
\centering
\small
\caption{Published constants and their Gomoku counterparts. The final column
marks whether a value is anchored in a primary source or is an experimental
choice that must be settled by measurement.}
\label{tab:scaling}
\begin{tabularx}{\textwidth}{@{}l l X@{}}
\toprule
Quantity & Gomoku value & Basis \\
\midrule
Simulations per move & 400 & experimental; half of AlphaZero's 800 \\
$c_{\mathrm{puct}}$ & 1.5 & not published by DeepMind; ELF OpenGo's value
\citep{tian2019elf}; tune by $\pm0.5$ \\
Dirichlet $\varepsilon$ & 0.25 & published \\
Dirichlet $\alpha$ & 0.1 & derived from $\alpha \approx 10/(\text{legal moves})$ \\
Temperature & $\tau=1$ for 12 moves, then $\to 0$ & derived by scaling AGZ's
30-of-200 to a 60-move game \\
Replay window & grow sublinearly from 250k positions & AlphaZero's window is
unpublished; adopt KataGo's growth rule \citep{wu2019katago} \\
Network & 128 channels, 10 residual blocks & derived; KataGo's second rung
\citep{wu2019katago}, well above AlphaGomoku's 32 filters $\times$ 2 blocks \\
Input planes & 2 to 4 & \code{me}, \code{you}, and optionally the last two moves \\
Gating & none & AlphaZero removed it \\
\bottomrule
\end{tabularx}
\end{table}
